\section{Recommendations}
\label{sec:recommendations}

\subsection{Numbered recommendations}

{\small
\begin{longtable}{@{}>{\raggedright\arraybackslash}p{0.8cm}>{\raggedright\arraybackslash}p{11.1cm}>{\raggedright\arraybackslash}p{2.2cm}>{\raggedright\arraybackslash}p{1.6cm}@{}}
\caption{Recommendations, owners (ERT teams, see Nomenclature) and timing.}
\label{tab:reco} \\
\toprule
\# & Recommendation & Owner & When \\
\midrule
\endfirsthead
\toprule
\# & Recommendation & Owner & When \\
\midrule
\endhead
R1 & Adopt the Von K\'arm\'an profile (Haack $C=0$). Use \texttt{nose\_profile\_final.csv} or the
STEP/STL solids (Section~\ref{sec:definition}) as the master geometry. & FD $\rightarrow$ ST & PDR \\
R2 & Book a \SI{907}{mm} nose bay (\SI{810}{mm} curved + \SI{97}{mm} cylinder for the 3U, of which
\SI{60}{mm} is the coupler shoulder) in the \SI{5.4}{m} LV budget. If it does not fit, keep the same
profile at $L = \SI{720}{mm}$ ($-\SI{60}{m}$, bay \SI{845}{mm}). & SE & PDR \\
R3 & Verify with Fluent (RANS, k-$\omega$ SST) at $M = 0.95$, 1.2 and 1.47 for the selected nose and
the power-$\tfrac12$ alternative, plus one $5^\circ$ incidence case for $C_{N\alpha}$ and CP, to
replace the transonic model of this report. & FD-CFD & before CDR \\
R4 & Put the final geometry and $C_{D,\rm w}(M)$ into the RocketPy model. Re-check apogee, rail-exit
speed and the 2--4 cal static margin once the fins are designed. The nose contributes
$C_{N\alpha} = 2$ at \FINXCP{} from the tip. & FD-Sim & PDR \\
R5 & Size the payload structure around a 3U envelope whose front face is \FINTHREEU{} ahead of the
shoulder, inside the $r \leq \SI{77}{mm}$ contour (Figure~\ref{fig:final}). & PL, ST & CDR \\
R6 & Create the missing Weisshorn nosecone requirements: profile, bay length and tip radius
(Section~\ref{sec:next}). & SE & PDR \\
\bottomrule
\end{longtable}
}

\subsection{When to choose another design}

The recommendation assumes the current Weisshorn priorities: apogee margin, a 3U payload in the
nose bay, and a length cap. If they change, Table~\ref{tab:alternatives} gives the design to use.
Do not use a cone or a tangent ogive on a supersonic Weisshorn, and judge a longer nose by the nose-bay length it needs, which is what the \SI{5.4}{m} cap constrains.

\begin{table}[H]
\centering
\caption{Choice of nose cone if the priorities change.}
\label{tab:alternatives}
\small
\begin{tabular}{@{}>{\raggedright\arraybackslash}p{6.3cm}>{\raggedright\arraybackslash}p{3.6cm}>{\raggedright\arraybackslash}p{5.3cm}@{}}
\toprule
If \dots & choose & cost relative to the recommendation \\
\midrule
the LV cannot grow beyond the LVS length & Von K\'arm\'an, \SI{720}{mm} & $-\SI{60}{m}$ apogee \\
the payload grows (more volume in the nose) & LV-Haack, \SI{810}{mm} & $-\SI{69}{m}$, bay $-\SI{37}{mm}$, $+\SI{22}{g}$ \\
the nose must be as short as possible ($L/D \leq 3.5$) & Power $\tfrac12$ & lowest drag at short length \\
the vehicle becomes faster ($M_{\max} \approx 2$) and the payload leaves the nose & Power $\tfrac34$ or Von K\'arm\'an at $L/D \geq 5$ & re-run the study \\
a subsonic test vehicle (e.g.\ Aletsch) & Tangent ogive or Von K\'arm\'an & shape barely matters below $M = 0.8$ \\
\bottomrule
\end{tabular}
\end{table}
